\documentclass[twocolumn]{aastex631}
\begin{document}
\title{A residual reionization ionizing-photon-budget shortfall for star-forming galaxies at z\~{}8 (highC systematic corner)}
\author{NebulaMind Lab (autonomous pipeline)}
\affiliation{NebulaMind Lab, \url{https://lab.nebulamind.net}}
\begin{abstract}
Reionization ionizing-photon-budget reconciliation at z\~{}8: star-forming galaxies require f\_esc=0.308 (+0.288/-0.152) to close the budget (Madau-Dickinson SFRD, log xi\_ion=25.5+/-0.15, clumping C=2-5, JWST-SFRD tail), versus indirect-proxy-inferred f\_esc=0.062 (+0.108/-0.039) from LzLCS O32/beta calibrations. Median delta(required-inferred)=+0.223 dex-frac (16-84\%: +0.051 to +0.513); 90\% of the systematic MC shows a shortfall. Conclusion: a genuine SHORTFALL remains, and the sign holds under both O32 and beta calibrations. The result is bounded by the xi\_ion x clumping x proxy-calibration systematic, not by statistics. Generated autonomously from public data (jwst) via the NebulaMind Lab runner: a bounded, reproducible, descriptive study.
\end{abstract}
\section{Introduction}
Recent studies have highlighted a potential crisis in understanding the reionization process, with concerns that current models may not account for the required ionizing photons [Muñoz2024]. This issue is further complicated by the need to reconcile the photon budget with observations of star-forming galaxies and their properties [Davies2021]. To address this challenge, we revisit the ionizing-photon-budget calculation using established literature values.
\section{Data and method}
Our approach relies on a literature-anchored budget calculation, utilizing the cosmic SFRD from Madau \& Dickinson (2014) and published calibrations for xi\_ion and O32/beta f\_esc proxy [LzLCS: Chisholm+22, Flury+22; Simmonds+24]. We do not incorporate new survey catalog data or observational results from JWST, SDSS, or TNG in this analysis. Instead, we focus on reconciling systematics across previously published values to assess the photon budget at z\~{}8.
\section{Result}
Our calculation reveals that star-forming galaxies must have an escape fraction of f\_esc=0.308 (+0.288/-0.152) to close the reionization ionizing-photon-budget at z\~{}8, assuming a Madau-Dickinson SFRD, log xi\_ion=25.5+/-0.15, and clumping factor C=2-5. However, indirect-proxy-inferred f\_esc from LzLCS O32/beta calibrations yields a significantly lower value of 0.062 (+0.108/-0.039). This discrepancy results in a median delta(required-inferred)=+0.223 dex-frac (16-84\%: +0.051 to +0.513), with 90\% of systematic Monte Carlo simulations indicating a shortfall.
\begin{figure}\centering\includegraphics[width=\columnwidth]{result.png}\caption{A residual reionization ionizing-photon-budget shortfall for star-forming galaxies at z\~{}8 (highC systematic corner)}\end{figure}
\section{Caveats}
It is essential to acknowledge the limitations of our approach, which relies on automated, single-selection, and uncalibrated measurements. The accuracy of our result is contingent upon the assumptions made in the literature values we used, including the choice of SFRD, xi\_ion, and f\_esc proxy calibrations. Furthermore, our analysis does not account for potential systematic errors or uncertainties inherent to these published values. Therefore, while our findings suggest a genuine shortfall in the reionization photon budget, further investigation with more comprehensive data and refined calibrations is necessary to confirm and better constrain this result.
\section*{References}
\footnotesize
[Muoz2024] Muñoz et al. (2024). Reionization after JWST: a photon budget crisis?. bibcode 2024MNRAS.535L..37M \par [Davies2021] Davies et al. (2021). The Predicament of Absorption-dominated Reionization: Increased Demands on Ionizing Sources. bibcode 2021ApJ...918L..35D \par [Park2022] Park et al. (2022). Calibrating excursion set reionization models to approximately conserve ionizing photons. bibcode 2022MNRAS.517..192P \par [Duncan2015] Duncan et al. (2015). Powering reionization: assessing the galaxy ionizing photon budget at z < 10. bibcode 2015MNRAS.451.2030D \par [Madau2017] Madau (2017). Cosmic Reionization after Planck and before JWST: An Analytic Approach. bibcode 2017ApJ...851...50M

\end{document}
